\documentclass[11pt, a4paper]{exam}
\usepackage[a4paper,top=2.3cm,bottom=2.8cm,left=2cm,right=2cm,headheight=15pt,headsep=0.6cm,footskip=1.2cm]{geometry}
\usepackage[T1]{fontenc}
\usepackage[utf8]{inputenc}
\usepackage{lmodern}
\usepackage{textcomp}
\usepackage[french]{babel}
\usepackage{xcolor}
\usepackage{listings}
\usepackage{graphicx}
\usepackage{array}
\usepackage{enumitem}
\usepackage{amsmath}
\usepackage{tikz}
\usetikzlibrary{positioning,shapes.multipart,arrows.meta}
\usepackage[most]{tcolorbox}

\definecolor{codegreen}{rgb}{0,0.6,0}
\definecolor{codegray}{rgb}{0.5,0.5,0.5}
\definecolor{codepurple}{rgb}{0.58,0,0.82}
\definecolor{backcolour}{rgb}{0.95,0.95,0.92}

\lstdefinestyle{mystyle}{
    backgroundcolor=\color{backcolour},
    commentstyle=\color{codegreen},
    keywordstyle=\color{magenta},
    numberstyle=\tiny\color{codegray},
    stringstyle=\color{codepurple},
    basicstyle=\ttfamily\footnotesize,
    breakatwhitespace=false,
    breaklines=true,
    postbreak=\mbox{\textcolor{codegray}{$\hookrightarrow$}\space},
    captionpos=b,
    keepspaces=true,
    columns=fullflexible,
    numbers=left,
    numbersep=5pt,
    showspaces=false,
    showstringspaces=false,
    showtabs=false,
    upquote=true,
    tabsize=4,
    xleftmargin=1.5em,
    framexleftmargin=0.3em,
    aboveskip=0.6em,
    belowskip=0.6em,
    extendedchars=true,
    literate={é}{{\'e}}1 {è}{{\`e}}1 {ê}{{\^e}}1 {ë}{{\"e}}1 {à}{{\`a}}1 {â}{{\^a}}1
             {ç}{{\c c}}1 {ù}{{\`u}}1 {û}{{\^u}}1 {ô}{{\^o}}1 {î}{{\^i}}1 {ï}{{\"i}}1
             {É}{{\'E}}1 {È}{{\`E}}1 {À}{{\`A}}1 {Ç}{{\c C}}1 {œ}{{\oe}}1
             {«}{{\guillemotleft}}1 {»}{{\guillemotright}}1 {→}{{$\rightarrow$}}1
}
\lstset{style=mystyle, language=Python}
\lstdefinestyle{sqlstyle}{style=mystyle, language=SQL, morekeywords={LIMIT,REFERENCES,IN,TEXT,FLOAT,DATE,VARCHAR,CHAR,REPLACE}, numbers=none, xleftmargin=0.3em, frame=lines}

% Code en ligne
\newcommand{\code}[1]{\texttt{#1}}

% Titre d'une question : numéro et intitulé, barème à droite
\newcommand{\Qt}[2]{\subsection*{#1\hfill{\normalfont\small #2}}}
% Rappel de la consigne en une ligne
\newcommand{\rappel}[1]{\noindent\textit{Ce qui est demandé :} #1\par\smallskip}

% La réponse attendue, mise en évidence
\newtcolorbox{reponse}[1][Réponse]{enhanced, breakable, colback=black!5, colframe=black, boxrule=0.9pt, arc=1mm,
  title={#1}, fonttitle=\bfseries\small, coltitle=black, colbacktitle=black!15, left=2mm, right=2mm, top=1mm, bottom=1mm,
  before skip=0.5em, after skip=0.7em}
% Erreur fréquente
\newtcolorbox{erreur}{enhanced, breakable, colback=white, colframe=black, boxrule=0.5pt, leftrule=3pt, arc=0mm,
  title={Erreur fréquente}, fonttitle=\bfseries\small, coltitle=black, colbacktitle=white, left=2mm, right=2mm, top=1mm, bottom=1mm,
  before skip=0.5em, after skip=0.7em}
% Encadré « À retenir »
\newtcolorbox{retenir}{enhanced, breakable, colback=white, colframe=black, boxrule=0.5pt, arc=1mm, boxsep=1mm,
  borderline={0.4pt}{1.2pt}{black, dashed},
  title={À retenir}, fonttitle=\bfseries\small, coltitle=black, colbacktitle=white, left=2mm, right=2mm, top=1mm, bottom=1mm,
  before skip=0.5em, after skip=0.7em}

\newcommand{\methode}{\noindent\textbf{Méthode.}\ }
\newcommand{\explication}{\noindent\textbf{Explication.}\ }

\setlist{itemsep=0.15em, topsep=0.3em}
\renewcommand{\arraystretch}{1.25}
\setlength{\parindent}{0pt}
\setlength{\parskip}{0.35em}

% Piles dessinées verticalement (sommet en haut) ; \suppr barre un élément supprimé
\newcommand{\suppr}[1]{\ooalign{#1\cr\hidewidth\rule[0.55ex]{1em}{1pt}\hidewidth\cr}}
\newcommand{\pilev}[2]{\begin{tikzpicture}[baseline=(titre.base)]
  \node[font=\scriptsize, align=center, text width=22mm] (titre) at (0,0) {#1};
  \foreach \v [count=\k] in {#2} {\node[draw, minimum width=8mm, minimum height=5mm, inner sep=0pt, font=\small] at (0,-0.25-\k*0.5) {\v};}
\end{tikzpicture}}
\setlength{\emergencystretch}{3em}
% Remarque sur l'énoncé (erreur ou ambiguïté du sujet)
\newtcolorbox{remarque}{enhanced, breakable, colback=white, colframe=black, boxrule=0.5pt, arc=0mm,
  borderline={0.4pt}{1.2pt}{black, dotted},
  title={Remarque sur l'énoncé}, fonttitle=\bfseries\small, coltitle=black, colbacktitle=black!8, left=2mm, right=2mm, top=1mm, bottom=1mm,
  before skip=0.5em, after skip=0.7em}
\newcommand{\sql}[1]{\textsc{#1}}
\lstset{style=sqlstyle}
\newcommand{\resultat}[1]{\par\smallskip\noindent\textit{Résultat sur le jeu d'essai :}\par\nopagebreak\begin{center}\small #1\end{center}}

\title{Corrigé — CC Term NSI - SQL : Gestion d'une Bibliothèque}
\author{Babash - NSI Tani Malandi et Johanne Bakalara}
\date{Contrôle de février 2026 (1 h)}

\begin{document}
\pagestyle{headandfoot}
\firstpageheader{}{}{}
\runningheader{\small Corrigé}{}{\small SQL, gestion d'une bibliothèque -- février 2026}
\runningheadrule
\firstpagefooter{}{\footnotesize Babash - NSI Tani Malandi et Johanne Bakalara \textperiodcentered{} CC BY-SA 4.0 \textperiodcentered{} page \thepage}{}
\runningfooter{}{\footnotesize Babash - NSI Tani Malandi et Johanne Bakalara \textperiodcentered{} CC BY-SA 4.0 \textperiodcentered{} page \thepage}{}

\maketitle
\thispagestyle{headandfoot}

\begin{tcolorbox}[colback=white, colframe=black, boxrule=0.5pt, arc=1mm, left=2mm, right=2mm]
\small\textbf{Comment lire ce corrigé.} Pour chaque question : ce qui est demandé, puis la réponse attendue (encadrée), puis l'explication pas à pas. Essaie d'abord seul·e, puis compare ta démarche, pas seulement ton résultat.
Toutes les requêtes ont été exécutées avec SQLite. Le sujet ne donne pas le contenu des tables : pour montrer un résultat, on utilise un petit \textbf{jeu d'essai} inventé (page suivante).
\end{tcolorbox}

\begin{remarque}
Le sujet annonce 7 points pour la partie 1 et 14 points pour la partie 2, soit \textbf{21 points} en tout. Ce corrigé reprend les points donnés par le sujet ; c'est à l'enseignant·e de ramener la note sur 20 s'il le faut.
\end{remarque}

\begin{retenir}
\begin{itemize}
\item Le schéma : Auteur(\underline{id\_auteur}, nom, prenom, nationalite) ; Livre(\underline{id\_livre}, titre, genre, prix, \#id\_auteur) ; Emprunt(\underline{id\_emprunt}, date\_emprunt, nom\_emprunteur, \#id\_livre).
\item Lire : \sql{SELECT} \ldots{} \sql{FROM} \ldots{} \sql{WHERE} \ldots{} \sql{ORDER BY} \ldots{} \sql{LIMIT} \ldots{} ;
\item Ajouter une ligne : \sql{INSERT INTO} ; modifier : \sql{UPDATE} \ldots{} \sql{SET} ; supprimer des lignes : \sql{DELETE FROM} ; supprimer la table entière : \sql{DROP TABLE}.
\end{itemize}
\end{retenir}

\section*{Partie 1 -- QCM (7 points)}

Pour chaque question, la bonne réponse est encadrée, puis on explique pourquoi les autres propositions sont fausses. Les lettres A, B, C désignent les propositions dans l'ordre du sujet.

\Qt{Question 1 -- Ajouter l'auteur Victor Hugo}{1 point}
\begin{reponse}
\textbf{A.} \code{INSERT INTO Auteur VALUES (1, 'Hugo', 'Victor', 'France');}
\end{reponse}
\explication \sql{INSERT INTO} ajoute une nouvelle ligne ; les valeurs suivent l'ordre des colonnes (id\_auteur, nom, prenom, nationalite). B (\sql{UPDATE}) modifie des lignes qui existent déjà : il ne crée pas d'auteur. C n'est pas du SQL valide : \sql{CREATE} sert à créer une table (\sql{CREATE TABLE}), pas à ajouter une ligne ; SQLite répond « syntax error ».

\Qt{Question 2 -- Supprimer l'emprunt numéro 102}{1 point}
\begin{reponse}
\textbf{B.} \code{DELETE FROM Emprunt WHERE id\_emprunt = 102;}
\end{reponse}
\explication \sql{DELETE FROM} supprime les lignes qui vérifient la condition du \sql{WHERE}. A et C provoquent une erreur de syntaxe : \sql{DROP} ne s'utilise pas avec \sql{FROM}, et \sql{UPDATE} ne supprime pas de ligne.

\begin{erreur}
Oublier le \sql{WHERE} : \code{DELETE FROM Emprunt;} supprime \textbf{tous} les emprunts (la table reste, mais vide).
\end{erreur}

\Qt{Question 3 -- Supprimer la structure de la table \code{Emprunt}}{1 point}
\begin{reponse}
\textbf{C.} \code{DROP TABLE Emprunt;}
\end{reponse}
\explication \sql{DROP TABLE} supprime la table elle-même (sa structure et toutes ses lignes) : après, \code{SELECT * FROM Emprunt;} donne l'erreur « no such table: Emprunt » (« table introuvable »). \code{DELETE TABLE} et \code{REMOVE TABLE} n'existent pas en SQL.

\Qt{Question 4 -- Passer le prix du livre 50 à 12.50}{1 point}
\begin{reponse}
\textbf{A.} \code{UPDATE Livre SET prix = 12.50 WHERE id\_livre = 50;}
\end{reponse}
\explication On modifie une ligne existante, donc \sql{UPDATE} \textit{table} \sql{SET} \textit{colonne} = \textit{valeur} \sql{WHERE} \textit{condition}. B est faux : \sql{INSERT} ajoute une ligne et n'accepte pas de \sql{WHERE}. C est faux : une requête ne peut pas commencer par \sql{SET}. En SQL, les nombres décimaux s'écrivent avec un point : \code{12.50}.

\Qt{Question 5 -- Les livres de genre 'BD' ou 'Manga'}{1 point}
\begin{reponse}
\textbf{A.} \code{WHERE genre IN ('BD', 'Manga');}
\end{reponse}
\explication \sql{IN} teste si la valeur appartient à une liste : c'est la même chose que \code{WHERE genre = 'BD' OR genre = 'Manga'}. B est faux : avec \sql{AND}, un livre devrait être à la fois BD \textbf{et} Manga, ce qui est impossible ; la requête ne renvoie aucune ligne. C n'est pas une syntaxe valide.

\begin{erreur}
Traduire « ou » du français par \sql{AND} parce qu'on veut « les BD et les mangas ». Pense ligne par ligne : chaque ligne doit être BD \textbf{ou} Manga.
\end{erreur}

\Qt{Question 6 -- La moyenne des prix}{1 point}
\begin{reponse}
\textbf{B.} \code{AVG(prix)}
\end{reponse}
\explication \sql{AVG} vient de l'anglais \textit{average}, « moyenne ». \sql{SUM} calcule la somme des prix et \sql{COUNT} compte les lignes. Sur le jeu d'essai : \code{SUM(prix)} vaut 164,2 ; \code{COUNT(prix)} vaut 7 ; \code{AVG(prix)} vaut 164,2 / 7, environ 23,46.

\Qt{Question 7 -- Les 3 livres les moins chers}{1 point}
\begin{reponse}
\textbf{A.} \code{ORDER BY prix ASC LIMIT 3;}
\end{reponse}
\explication On trie du moins cher au plus cher (\sql{ASC}, \textit{ascending}, « croissant »), puis on garde les 3 premières lignes avec \sql{LIMIT} 3. B trie par prix décroissant (\sql{DESC}) : on obtiendrait les 3 livres les \textbf{plus} chers. C affiche le nombre 3 sur chaque ligne de la table : \code{SELECT 3} ne limite rien.

\section*{Partie 2 -- Rédaction de requêtes (14 points)}

\textbf{Jeu d'essai inventé} (les tables \code{Auteur}, \code{Livre} et \code{Emprunt} remplies pour tester les requêtes) :
\begin{center}\footnotesize
\begin{tabular}{|c|l|l|l|}\hline
\multicolumn{4}{|c|}{\textbf{Auteur}}\\\hline
id\_auteur & nom & prenom & nationalite \\\hline
1 & Hugo & Victor & France \\
2 & Madi & Zaïnaba & France \\
3 & Sato & Ren & Japon \\
4 & Okafor & Ada & Nigeria \\\hline
\end{tabular}\quad
\begin{tabular}{|c|l|l|l|}\hline
\multicolumn{4}{|c|}{\textbf{Emprunt}}\\\hline
id\_emprunt & date\_emprunt & nom\_emprunteur & id\_livre \\\hline
101 & 2026-02-02 & Anissa & 12 \\
102 & 2026-02-03 & Yanis & 10 \\
103 & 2026-02-05 & Anissa & 13 \\
104 & 2026-02-06 & Kenza & 11 \\\hline
\end{tabular}\\[0.6em]
\begin{tabular}{|c|l|l|r|c|}\hline
\multicolumn{5}{|c|}{\textbf{Livre}}\\\hline
id\_livre & titre & genre & prix & id\_auteur \\\hline
10 & Les Misérables & Roman & 25.0 & 1 \\
11 & Le Secret de la falaise & Roman & 18.5 & 2 \\
12 & Lagon Quest, tome 1 & Manga & 7.2 & 3 \\
13 & La Route du sud & Roman & 22.0 & 4 \\
14 & Le secret des étoiles & BD & 14.0 & 2 \\
15 & Notre-Dame de Paris, édition illustrée & Roman & 65.0 & 1 \\
50 & Chronique des sept misères & Roman & 12.5 & 2 \\\hline
\end{tabular}
\end{center}
{\small Les auteur·rices 2 à 4, les titres 11 à 14 et 50 ainsi que les emprunteur·euses sont fictifs.}

\Qt{Question 8 -- Créer la table \code{Emprunt}}{2 points}
\rappel{créer la table avec ses types, sa clé primaire et sa clé étrangère.}
\begin{reponse}
\begin{lstlisting}
CREATE TABLE Emprunt (
    id_emprunt INT PRIMARY KEY,
    date_emprunt TEXT,
    nom_emprunteur TEXT,
    id_livre INT,
    FOREIGN KEY (id_livre) REFERENCES Livre(id_livre)
);
\end{lstlisting}
\end{reponse}
\explication
\begin{itemize}
\item On écrit chaque attribut suivi de son type, comme sur le schéma (INT, TEXT), séparés par des virgules.
\item \sql{PRIMARY KEY} indique la clé primaire (\code{id\_emprunt}, soulignée dans le schéma).
\item \sql{FOREIGN KEY} (\code{id\_livre}) \sql{REFERENCES} \code{Livre(id\_livre)} indique la clé étrangère (\#) : un emprunt doit porter sur un livre qui existe. Avec cette contrainte, SQLite refuse d'enregistrer un emprunt du livre 999, qui n'existe pas (« FOREIGN KEY constraint failed »).
\item La table \code{Livre} doit exister avant \code{Emprunt}, puisque \code{Emprunt} y fait référence.
\end{itemize}

\begin{erreur}
Mettre une virgule après la dernière ligne, juste avant la parenthèse fermante : c'est une erreur de syntaxe. Autre oubli fréquent : le point-virgule final, exigé par les consignes.
\end{erreur}

\Qt{Question 9 -- Les titres de tous les livres}{2 points}
\begin{reponse}
\begin{lstlisting}
SELECT titre FROM Livre;
\end{lstlisting}
\end{reponse}
\explication On ne veut qu'une colonne (\code{titre}) et toutes les lignes : pas de \sql{WHERE}. \code{SELECT *} afficherait toutes les colonnes, alors que la consigne dit « uniquement les titres ». Sur le jeu d'essai, la requête donne les 7 titres.

\Qt{Question 10 -- Les romans à plus de 20 \texteuro}{2 points}
\begin{reponse}
\begin{lstlisting}
SELECT titre FROM Livre WHERE genre = 'Roman' AND prix > 20;
\end{lstlisting}
\end{reponse}
\resultat{\begin{tabular}{|l|}\hline titre \\\hline Les Misérables \\ La Route du sud \\ Notre-Dame de Paris, édition illustrée \\\hline\end{tabular}}
\explication Les deux conditions doivent être vraies en même temps : on les relie par \sql{AND}. « Strictement supérieur » s'écrit \code{>} ; \code{>=} garderait aussi les livres à 20 \texteuro{} pile. Le texte \code{'Roman'} est entre guillemets, le nombre \code{20} non (et sans le symbole \texteuro).

\Qt{Question 11 -- Le nombre d'enregistrements de \code{Emprunt}}{2 points}
\begin{reponse}
\begin{lstlisting}
SELECT COUNT(*) FROM Emprunt;
\end{lstlisting}
\end{reponse}
\resultat{\begin{tabular}{|c|}\hline COUNT(*) \\\hline 4 \\\hline\end{tabular}}
\explication \code{COUNT(*)} compte les lignes de la table. La requête renvoie un seul nombre, pas la liste des emprunts.

\Qt{Question 12 -- Livres à plus de 50 \texteuro{} et leur auteur}{2 points}
\begin{reponse}
\begin{lstlisting}
SELECT Livre.titre, Auteur.nom
FROM Livre
JOIN Auteur ON Livre.id_auteur = Auteur.id_auteur
WHERE Livre.prix > 50;
\end{lstlisting}
\end{reponse}
\resultat{\begin{tabular}{|l|l|}\hline titre & nom \\\hline Notre-Dame de Paris, édition illustrée & Hugo \\\hline\end{tabular}}
\methode
\begin{enumerate}
\item Le titre est dans \code{Livre}, le nom de l'auteur dans \code{Auteur} : il faut une \textbf{jointure}.
\item On relie les deux tables par la clé étrangère \code{Livre.id\_auteur} et la clé primaire \code{Auteur.id\_auteur}.
\item On garde les livres de prix supérieur à 50.
\end{enumerate}
La consigne dit « supérieur à 50 » sans « strictement » : \code{prix >= 50} est aussi accepté.

\begin{erreur}
Écrire \code{SELECT titre, nom FROM Livre, Auteur WHERE prix > 50;} sans condition de jointure : chaque livre est associé à \textbf{chaque} auteur, et Notre-Dame de Paris apparaîtrait 4 fois, avec 4 noms d'auteur·rices différents.
\end{erreur}

\Qt{Question 13 -- Les titres qui contiennent « Secret »}{2 points}
\begin{reponse}
\begin{lstlisting}
SELECT titre FROM Livre WHERE titre LIKE '%Secret%';
\end{lstlisting}
\end{reponse}
\resultat{\begin{tabular}{|l|}\hline titre \\\hline Le Secret de la falaise \\ Le secret des étoiles \\\hline\end{tabular}}
\explication \sql{LIKE} compare à un motif ; \code{\%} remplace n'importe quelle suite de caractères. \code{'\%Secret\%'} veut dire : n'importe quoi, puis « Secret », puis n'importe quoi, donc « contient Secret ».
Remarque : SQLite ne distingue pas les majuscules dans \sql{LIKE}, c'est pourquoi « secret » est aussi trouvé ; d'autres logiciels ne trouveraient que « Le Secret de la falaise ».

\begin{erreur}
Écrire \code{WHERE titre = 'Secret'} : on ne trouve que les livres dont le titre est exactement « Secret ». Ou oublier un des deux \code{\%} : \code{'Secret\%'} cherche les titres qui \textbf{commencent} par « Secret ».
\end{erreur}

\Qt{Question 14 -- Tous les livres triés par titre}{2 points}
\begin{reponse}
\begin{lstlisting}
SELECT * FROM Livre ORDER BY titre ASC;
\end{lstlisting}
\end{reponse}
\explication \code{*} affiche toutes les informations (toutes les colonnes). \sql{ORDER BY} \code{titre} trie les lignes selon le titre ; \sql{ASC} (ordre croissant, de A à Z) est l'ordre par défaut et peut être omis. Sur le jeu d'essai, l'ordre obtenu est : Chronique des sept misères, La Route du sud, Lagon Quest, Le Secret de la falaise, Le secret des étoiles, Les Misérables, Notre-Dame de Paris.

\begin{erreur}
Placer \sql{ORDER BY} avant \sql{FROM} ou avant \sql{WHERE} : l'ordre des mots est imposé, \sql{SELECT} \ldots{} \sql{FROM} \ldots{} \sql{WHERE} \ldots{} \sql{ORDER BY} \ldots
\end{erreur}

\end{document}
